% ADASS 2026 contributed talk C803.
% Based on the ADASS_template.tex LaTeX file (24 Oct 2025) and the ASP 2021 author template.

\documentclass[11pt,twoside]{article}

% Do NOT use ANY packages other than asp2023.
\usepackage{asp2023}

\aspSuppressVolSlug
\resetcounters

\bibliographystyle{asp2021}

\markboth{Chattopadhyay}{FIBRA: Reduction of Fiber-fed IFU Spectra}

\begin{document}

\title{FIBRA: FIber Bundle Reduction of Astronomical spectra --- an Instrument-agnostic Pipeline for Fiber-fed Integral-field Spectroscopy}

\author{Sabyasachi~Chattopadhyay}
\affil{University of Hawai`i at M\=anoa, Honolulu, HI, USA; \email{sabya@hawaii.edu}}

\paperauthor{Sabyasachi~Chattopadhyay}{sabya@hawaii.edu}{}{University of Hawaii at Manoa}{Space Science and Engineering Initiative}{Honolulu}{HI}{96822}{USA}

%\aindex{Chattopadhyay,~S.}

\begin{abstract}
FIBRA is a Python package for the reduction of data from fiber-fed integral field units (IFUs). An instrument enters as data: a detector description, a fiber map and rules for typing frames from their headers. Twelve stages lead from the raw frame to a flux-calibrated data cube. The variance starts at the detector and is propagated through every stage, and each product carries a cumulative key of its inputs and parameters. Every interactive stage has a tuner that writes the same configuration file the batch run reads. Three methods address effects met with the SMI-300 IFU on the Robert Stobie Spectrograph of the Southern African Large Telescope: lamp and sky illuminate the focal plane differently, the sky varies across the field, and neighboring fibers overlap on a dense pseudo-slit. We compare each stage with IRAF on identical input, and the final cube with the MaNGA cube of the same galaxy.
\end{abstract}

%\ssindex{data!pipelines}
%\ssindex{algorithm!data reduction}
%\ssindex{data!provenance}
%\ssindex{spectroscopy!integral field}
%\ssindex{instruments!Integral Field Unit (IFU)}
%\ssindex{techniques!sky subtraction}
%\ssindex{observatories!ground-based!Southern African Large Telescope (SALT)}
%\ssindex{packages!IRAF}

%\ooindex{IRAF, ascl:9911.002}
%\ooindex{R3D, ascl:1106.005}
%\ooindex{p3d, ascl:1205.002}
%\ooindex{L.A.Cosmic, ascl:1207.005}
%\ooindex{PyCosmic, ascl:1210.027}
%\ooindex{Astropy, ascl:1304.002}

\section{Introduction}

Fiber-fed IFUs reformat a field into one or more pseudo-slits at the entrance of a spectrograph. All fiber-fed spectrographs yield similar raw data, and their reduction follows a similar sequence of steps \citep{2006AN....327..850S}. Several packages serve more than one instrument, among them the IRAF \texttt{dohydra} tasks \citep{1993ASPC...52..173T}, R3D \citep{2006AN....327..850S}, p3d \citep{2010A&A...515A..35S} and PyKOALA \citep{2025arXiv250718347C}. Survey instruments carry dedicated pipelines, such as those of MaNGA \citep{2016AJ....152...83L} and SAMI \citep{2015MNRAS.446.1551S}.

Fiber IFUs added to general-purpose spectrographs share the slit plane, the calibration unit and the detector of the modes they join. Three effects can follow from this arrangement. First, the calibration lamp can illuminate the focal plane differently from the sky. At the Southern African Large Telescope (SALT), the lamp passes a moving exit-pupil baffle and only approximates the sky illumination \citep{2008SPIE.7014E..6HB}. Second, sky fibers placed away from the object field can see a different sky. On one SMI-300 exposure, two sky bundles on either side of the field reach about 141 and 190 counts at 6000\,\AA. Third, a dense pseudo-slit can bring neighboring traces closer than specified, and a fixed aperture then trades lost light against cross-talk. In the documentation of the packages above we found partial support, or none, for these three effects. FIBRA was written so that a fiber IFU enters as an instrument description rather than as code.

\section{Design}

FIBRA has three layers. Pure stage functions take an input and parameters and return an output. One engine runs these functions, caches their products, records provenance and writes a quality report. Two front ends drive the same engine: a command-line interface and a graphical interface with eight tabs. The stage chain is detector processing, cosmic-ray removal, extraction, flat fielding, wavelength calibration, flexure, sky subtraction, standard-star extraction, flux calibration, telluric correction, cube building and frame combination.

One YAML file holds the configuration, with one block per stage. Its keys are the parameters shown on each stage's page of the interface, and values are type-checked on loading. The instrument description is data. It gives the detector (overscan, gain, read noise), the fiber map (trace order, fiber type, sky position, relative throughput, illumination factor), header rules for frame types, and defaults. SMI-200, SMI-300 and a synthetic test instrument run from such descriptions.

Every product is one of four containers: Image2D, RSS, Cube and Spectrum. Each carries flux, variance, a mask, a header and a history. The variance starts at detector processing as $\sigma^2 = \max(N_e, 0) + R^2$, with $N_e$ the counts in electrons and $R$ the read noise. It is then propagated through the flat field, the sky model, the sensitivity function and the resampling. The cube adds a covariance factor following \citet{2015MNRAS.446.1551S}, a per-spaxel line-spread function and a celestial WCS.

Each product also carries a cumulative SHA-1 key of the input frame, the ordered stage parameters and a code revision. A rerun reuses every product whose key is unchanged. Changing one sky-subtraction parameter therefore reprocesses the chain from sky subtraction onward, and the six earlier products are reused.

Seven stages have tuner windows, from cosmic rays to the sensitivity function. A value accepted in a tuner is written to the configuration file, so the batch run reproduces the interactive session.

\section{Methods}

\textbf{Focal-plane illumination.} The extracted flat of fiber $i$ is multiplied by $\Gamma_i = T_i / L_i$ before its continuum is normalized. Here $T_i$ is the relative sky illumination at the fiber's position on the pseudo-slit, and $L_i$ the relative lamp illumination. For the SMIs, $\Gamma_i$ comes from the ratio of a sky emission line between a long-slit exposure and an SMI exposure. Without a long-slit exposure, a flat and a sky RSS give $L_i$ at the central column and $T_i$ at the brightest airglow line. For SMI-300, $\Gamma_i$ spans 1.00 to 1.14. The relative fiber throughput is kept separate from $\Gamma_i$.

\textbf{Sky subtraction in field zones.} Object and sky fibers are grouped into zones by position in the field, with the sky fibers grouped by a friends-of-friends link on the sky \citep{1982ApJ...257..423H}. Each zone gets a clipped continuum model and a spatial model of every sky line across the slit rows. Each object fiber then takes its sky from the fibers of its own zone. The residuals are compared with the photon noise zone by zone. SMI-300 uses eight zones built from its 34 sky fibers.

\textbf{Variable apertures and aperture loss.} Extraction is a boxcar, optimal \citep{1986PASP...98..609H} or a joint fit over neighboring fibers \citep{2010PASA...27...91S}. Each fiber has its own lower and upper aperture bounds and center offset. These can vary at up to 50 wavelengths, and can be edited with the mouse or as a table. An aperture-loss correction is then applied to science and standard-star spectra alike. It uses normalized pinhole profiles of the fibers, the trace separations, and 98\,\% enclosed energy as the total extent of a fiber's light along the slit. On SMI-300, traces 243 and 244 sit 2\,px closer than specified, and H$\alpha$ from trace 243 enters the aperture of trace 244. Widening the boxcar from 4 to 6\,px raises the kept light of trace 244 from 69 to 88\,\% of its total. It also raises the leak from trace 243 from 25 to 32\,\% of its H$\alpha$.

\textbf{Other stages.} Cosmic rays are rejected on one frame with the Laplacian method of \citet{2001PASP..113.1420V} or across repeated frames. Scattered light is fitted with splines through the slit edges and the gaps left by missing fibers. Arc lines on a reference fiber are identified with the Hough transform and RANSAC approach of \citet{2019arXiv191205883V}, and the solution is then carried to every fiber. On SMI-300, 31 xenon lines give a median residual of 0.34\,\AA\ over the 244 fibers. Sky lines in each science frame measure the flexure shift. The standard star is extracted over the IFU with a \citet{1969A&A.....3..455M} profile whose center follows atmospheric refraction with wavelength. The same star provides the sensitivity function and the telluric correction. The cube is built with a modified-Shepard Gaussian kernel as in \citet{2016AJ....152...83L}, or with drizzle \citep{2002PASP..114..144F}.

\section{Tests}

The tests use the SMI-300 slit-mask IFU, which we built for the Robert Stobie Spectrograph \citep[RSS;][]{2003SPIE.4841.1463B} on SALT \citep{2024SPIE13096E..2RC, 2026arXiv260322698C}. Its fiber map lists 244 fibers: 210 on an object field of about $28\arcsec \times 17\arcsec$ and 34 in two sky bundles about $80\arcsec$ away. We compare each stage with IRAF 2.17 \citep{1993ASPC...52..173T} on an SMI-300 science frame and a standard-star frame. IRAF receives FIBRA's exact input for the stage under test, with matched parameters.

\begin{itemize}
\item Extraction (\texttt{apall}): the traces agree to 0.04\,px rms. The flat fluxes differ by $+0.065$\,\% through partial-pixel weighting at the aperture edges. The arc and standard spectra agree within $0.1\sigma$.
\item Flat field (the \texttt{msresp1d} response rebuilt with \texttt{sfit} and \texttt{imarith}): the responses agree within $\pm 0.03$\,\%, and the flat-fielded science spectra within $3\times10^{-3}\sigma$.
\item Wavelength (\texttt{autoidentify}, \texttt{reidentify}, \texttt{dispcor}, 31 lines): the solutions agree to 0.18\,\AA\ rms. The arc-fit rms is 0.42\,\AA\ for FIBRA and 0.50\,\AA\ for IRAF.
\item Sky (\texttt{skysub} with one sky, against FIBRA with one zone): the residual scatter at sky lines agrees within $0.02\sigma$ ($0.58\sigma$ and $0.60\sigma$).
\item Flux (\texttt{standard}, \texttt{sensfunc}, \texttt{calibrate}): the IRAF result is 1.3\,\% higher in the median, with a $\pm 6$\,\% ripple from its order-17 fit.
\item Telluric (\texttt{telluric}, same settings): the corrected spectra agree to $3\times10^{-7}$.
\end{itemize}

We built a cube of the MaNGA galaxy 12510-9102 from three dithered SMI-300 exposures and placed it beside the MaNGA DR17 cube \citep{2022ApJS..259...35A} over the same sky area. After a rotation of $180\deg$ and registration by cross-correlation, the H$\alpha$ maps show the same knots. The H$\alpha$ profiles at matched positions have the same shape, and the SMI-300 line is broader, in line with its lower resolving power. A quantitative comparison of fluxes and kinematics is the subject of a separate paper on SMI-300 science verification.

\section{Documentation and Status}

The Help tab of the interface has 35 pages. They cover the eight tabs, two dialogs, seven tuner windows, fifteen configuration sections, the command line and the product files. The lists of 212 parameters and 204 interface controls are generated from the code, and a completeness check lists any without a description. A 58-page beginner's guide walks through a first reduction, with window captures from the running code.

FIBRA has about 20\,000 lines of Python and depends on NumPy \citep{2020Natur.585..357H}, SciPy \citep{2020NatMe..17..261V}, Astropy \citep{2022ApJ...935..167A}, PyYAML and Matplotlib \citep{2007CSE.....9...90H}. The University of Hawai`i has cleared FIBRA for public release, and a release with a DOI and continuous integration is planned. Further fiber IFUs would enter as instrument descriptions.

\acknowledgements Based on observations made with the Southern African Large Telescope (SALT). This work uses data from SDSS-IV MaNGA DR17.

\bibliography{C803}

% if we have space left, we might add a conference photograph here. Leave commented for now.
% \bookpartphoto[width=1.0\textwidth]{foobar.eps}{FooBar Photo (Photo: Any Photographer)}

\end{document}
